%% Resume.tex — Résumé / Abstract
\chapter*{Résumé}
\addcontentsline{toc}{chapter}{Résumé}

\par Dans un environnement économique de plus en plus orienté vers la donnée, la
capacité d'une entreprise à analyser sa performance financière, à se comparer à
ses concurrents et à anticiper ses résultats constitue un avantage stratégique
déterminant. Les états financiers, riches en information, restent toutefois
difficiles à exploiter sous leur forme brute : volumineux, hétérogènes d'une
année à l'autre et peu propices à l'analyse comparative ou prédictive.

\par Ce projet de fin d'études, mené dans le cadre du Master professionnel en
\textbf{Business Analytics} d'\textbf{Esprit School of Business}, propose la
conception et la réalisation d'une plateforme décisionnelle d'analyse financière
pour la \textbf{Société Frigorifique et Brasserie de Tunis (SFBT)}. À partir des
états financiers de la société sur la période \textbf{2005--2025}, une chaîne de
traitement automatisée a été développée en \textit{Python}~: nettoyage et
unification de \textbf{quatre sources hétérogènes} (la SFBT et trois sociétés de
\textit{benchmark}~: DELICE, AB~InBev et Coca-Cola), fiabilisation des données
(correction des doublons et des colonnes comparatives erronées), puis
\textbf{augmentation} par génération des clôtures trimestrielles manquantes selon
une méthode distinguant les grandeurs de stock et de flux. Le cœur analytique
calcule \textbf{46 ratios financiers} conformes à un référentiel industriel
normalisé, assortis de fourchettes de référence et d'une grille de notation,
agrégés en un \textbf{Score Global de Performance Industrielle (SGPI)} sur 100.
Les données sont organisées dans un \textbf{entrepôt de données} en schéma en
étoile, interrogeable en SQL et destiné à alimenter des tableaux de bord
\textit{Power BI}. Enfin, un module de \textbf{prévision} compare trois modèles
d'apprentissage automatique (régression linéaire, forêt aléatoire et arbre de
décision) pour estimer les valeurs financières futures.

\par La solution, entièrement \textbf{testée} (83 tests automatiques, tous
satisfaits) et \textbf{documentée}, place la SFBT en tête de son secteur avec un
score de \textbf{83,5/100}, et atteint une précision de prévision de 85 à 96~\%.
Elle fournit ainsi un diagnostic financier chiffré, comparatif et prédictif de
la SFBT, et constitue une base reproductible et extensible pour le pilotage de
la performance.

\vspace{0.8cm}
{\large \textbf{Mots-clés :}} Business Intelligence, entrepôt de données, schéma
en étoile, ratios financiers, MSI20000, \textit{benchmark}, apprentissage
automatique, prévision, SFBT, Python, Power BI.

\newpage

%% ── Abstract en anglais ────────────────────────────────────────────────────
\selectlanguage{english}
\chapter*{Abstract}
\addcontentsline{toc}{chapter}{Abstract}

\par In an increasingly data-driven economic environment, a company's ability to
analyze its financial performance, benchmark itself against competitors and
anticipate its results is a decisive strategic advantage. Financial statements,
although rich in information, remain difficult to exploit in their raw form:
voluminous, heterogeneous from one year to the next, and ill-suited to
comparative or predictive analysis.

\par This end-of-study project, carried out within the professional Master's in
\textbf{Business Analytics} at \textbf{Esprit School of Business}, presents the
design and implementation of a financial-analysis decision-support platform for
the \textbf{Société Frigorifique et Brasserie de Tunis (SFBT)}. Based on the
company's financial statements over \textbf{2005--2025}, an automated
\textit{Python} pipeline was developed: cleaning and unification of \textbf{four
heterogeneous sources} (SFBT and three benchmark companies: DELICE, AB~InBev and
Coca-Cola), data reliability fixes (duplicate removal and correction of erroneous
comparative columns), then \textbf{data augmentation} through the generation of
the missing quarterly closings, using a method that distinguishes stock from flow
quantities. The analytical core computes \textbf{46 financial ratios} compliant
with a normalized industrial framework, together with reference ranges and a
scoring grid, aggregated into a \textbf{Global Industrial Performance Score} out
of 100. The data is organized in a \textbf{star-schema data warehouse}, queryable
in SQL and designed to feed \textit{Power BI} dashboards. Finally, a
\textbf{forecasting} module compares two machine-learning models (linear
regression and random forest) to estimate future financial values.

\par The solution, fully \textbf{tested} (70 automated tests, all passing) and
\textbf{documented}, ranks SFBT first in its sector with a score of
\textbf{83.5/100}, and reaches a forecasting accuracy of 85 to 96\%. It thus
delivers a quantified, comparative and predictive diagnosis of SFBT's financial
performance, providing a reproducible and extensible foundation for performance
management.

\vspace{0.8cm}
{\large \textbf{Keywords:}} Business Intelligence, data warehouse, star schema,
financial ratios, MSI20000, benchmark, machine learning, forecasting, SFBT,
Python, Power BI.

\selectlanguage{french}
\newpage
